\documentclass[11pt]{article}
\usepackage[margin=1in]{geometry}
\usepackage{amsmath,amssymb,amsthm,mathtools,bm}
\usepackage{enumitem,booktabs,graphicx,xcolor}
\usepackage[numbers,sort&compress]{natbib}
\usepackage[colorlinks=true,linkcolor=blue!50!black,citecolor=blue!50!black,urlcolor=blue!50!black]{hyperref}
\usepackage[capitalise]{cleveref}

\newtheorem{theorem}{Theorem}
\newtheorem{lemma}[theorem]{Lemma}
\newtheorem{proposition}[theorem]{Proposition}
\newtheorem{corollary}[theorem]{Corollary}
\newtheorem{conjecture}[theorem]{Conjecture}
\theoremstyle{definition}
\newtheorem{definition}[theorem]{Definition}
\newtheorem{assumption}[theorem]{Assumption}
\newtheorem{procedure}[theorem]{Procedure}
\theoremstyle{remark}
\newtheorem{remark}[theorem]{Remark}
\newtheorem{example}[theorem]{Example}
\crefname{procedure}{Procedure}{Procedures}
\crefname{assumption}{Assumption}{Assumptions}

\newcommand{\R}{\mathbb{R}}
\newcommand{\E}{\mathbb{E}}
\newcommand{\N}{\mathcal{N}}
\newcommand{\T}{\mathcal{T}}
\newcommand{\F}{\mathfrak{F}}
\newcommand{\cA}{\mathcal{A}}
\newcommand{\hy}{\hat{y}}
\newcommand{\diam}{\operatorname{diam}}
\newcommand{\tr}{\operatorname{tr}}
\newcommand{\Var}{\operatorname{Var}}

\title{Cobweb Forests:\\ A Kolmogorov--Arnold Theory of Networks of Concept Hierarchies}
\author{Draft --- Karthik Singaravadivelan et al.\ (working notes)}
\date{\today}

\begin{document}
\maketitle

\begin{abstract}
The Kolmogorov--Arnold (KA) representation theorem states that every continuous
function of $n$ variables is a finite superposition of continuous functions of one
variable and addition, and Kolmogorov--Arnold Networks (KANs) turn that statement into
an architecture with learnable univariate functions on its edges.  We ask the two
analogous questions for Cobweb, the incremental concept-formation system whose
category utility is the mutual information between features and concepts.
\emph{Representation:} is there a set of Cobweb trees that represents an arbitrary
continuous function?  \emph{Learnability:} does Cobweb's own learning
procedure---category utility and its four operators, with no gradients---learn it?
We formalise a \emph{Cobweb forest}: a directed acyclic graph of hybrid Cobweb trees in
which each tree reads raw features and the \emph{messages} of other trees---their
\emph{activations} (multi-node predictions of an attribute) and their \emph{labels}
(the concepts an input sorts to)---and one designated tree produces the output.
On representation we prove that a single tree is a universal approximator under two of
the implementation's four prediction weightings (under the other two the natural
construction keeps a fixed bias, and so do learned trees), a
Cobweb--Kolmogorov--Arnold theorem (forests of two-input trees passing activations are
universal with $4n^2+4n$ trees, and one input per tree is not enough), that label forests
are universal, and that merging trees under a common root yields mixtures, not
compositions.  On learnability we prove that no Cobweb operator ever deletes a leaf,
which makes a learned tree a universally consistent regressor whatever structure
category utility builds, and we introduce \emph{cognitive tasks}: what each tree
perceives (its \emph{cue}) and what it predicts and later observes (its
\emph{outcome}).  A forest learned by Cobweb converges to the \emph{ideal forest} of its
task assignment, which equals the target exactly when the assignment is
\emph{sufficient}; predicting the target at every tree is sufficient for
monotone-additive targets, predicting another view of the input is sufficient under
multi-view redundancy, and categorization is always sufficient but at a rate governed by
the dimension.  Tasks cannot be inferred locally---a sign construction makes a tree's
own data identical whether its input is necessary or irrelevant---but they can be
inferred consistently by verifying candidate tasks at the tree that composes them.  An
executable verifier checks each statement against the implementation.
\end{abstract}

\tableofcontents

%=====================================================================
\section{Introduction}
%=====================================================================

\paragraph{Kolmogorov--Arnold.}
For every $n\ge 2$ there are continuous functions $\varphi_{q,p}:[0,1]\to\R$
($q=0,\dots,2n$, $p=1,\dots,n$) such that every $f\in C([0,1]^n)$ can be written as
\begin{equation}\label{eq:ka}
  f(x_1,\dots,x_n)=\sum_{q=0}^{2n}\Phi_q\Bigl(\sum_{p=1}^{n}\varphi_{q,p}(x_p)\Bigr)
\end{equation}
with continuous outer functions $\Phi_q:\R\to\R$ that depend on~$f$
\citep{kolmogorov1957,arnold1957,lorentz1966,sprecher1965}.  The only genuinely
multivariate operation is addition.  The inner functions may be very rough (nowhere
differentiable), which long made \eqref{eq:ka} look irrelevant to learning
\citep{girosi1989}.  Kolmogorov--Arnold Networks \citep{liu2024kan} stack layers of
learnable univariate functions (B-splines on edges, sums at nodes) and prove an
approximation theorem whose rate does not depend on $n$ when $f$ has a \emph{smooth}
KA-type representation \citep[Thm.~2.1]{liu2024kan}.

\paragraph{Cobweb.}
Cobweb \citep{fisher1987} sorts each instance down a hierarchy of probabilistic
concepts and, at every node, applies whichever of four operators---add to the best
child, create a new child, merge the two best children, split the best child---most
increases category utility, which in its information-theoretic form is the mutual
information between features and concepts \citep{corter1992,maclellan2026rational}.
Concepts store sufficient statistics (attribute--value counts for nominal attributes,
a mean and variance for each continuous one), so learning is closed-form, incremental
and does not forget.  Prediction expands concepts best-first and mixes their
predictions \citep{maclellan2026rational}.

Several recent systems use Cobweb trees as \emph{modules} whose representations become
the inputs of other trees: Convolutional Cobweb \citep{maclellan2021conv} sorts image
patches into a filter hierarchy and describes each image by the patches' filter labels
for a second, classification hierarchy; TRELLIS \citep{singaravadivelan2026trellis}
describes composites by bags of concept identifiers drawn from another hierarchy; and a
Cobweb language model passes each tree's prediction for a position to the trees above
it.  We call such a network a \emph{Cobweb forest} and ask what it can represent and
learn.

\paragraph{Questions.}
\begin{enumerate}[label=\textbf{Q\arabic*.},leftmargin=*]
\item \emph{Representation.} For every continuous $f$ and $\varepsilon>0$, is there a
  Cobweb forest whose output is uniformly within $\varepsilon$ of $f$?  How few trees,
  and how few inputs per tree, suffice?
\item \emph{Learnability.} If every tree learns only with Cobweb's own procedure, does
  the forest's output converge to $f$?  What must each tree be asked to predict, and can
  that be inferred from data?
\end{enumerate}

\paragraph{Representation is not learnability.}
The constructions that answer Q1 choose what every tree holds knowing $f$: the trees of
the KA construction hold the values $\varphi_{q,p}(x_p)$ and $\Phi_q(u)$.  A learner never
observes these values; it observes its inputs and, eventually, outcomes.  A KAN recovers
the individual terms of a sum from the error of the sum by back-propagation.  A Cobweb
tree has no gradient: it can only learn to predict what it will observe.  Q2 is
therefore a question about \emph{tasks}---what each tree perceives and what it
predicts---and the central objects of the second half of the paper are \emph{cognitive
tasks} (\cref{sec:tasks}) and their inference (\cref{sec:infer}).

\paragraph{Contributions.}
\begin{itemize}[leftmargin=*]
\item A formal model of the hybrid Cobweb tree as implemented in
  \texttt{cobweb\_hybrid}, including its four prediction weightings, and of Cobweb
  forests with activation and label messages (\cref{sec:model}).  An executable verifier
  checks that the model equals the code (\cref{app:verify}).
\item Representation (\cref{sec:repr}): single-tree universality with an explicit
  rate (\cref{thm:single}); inconsistency of the prior-weighted all-concept mixture
  (\cref{prop:prior}); the Cobweb--Kolmogorov--Arnold theorem (\cref{thm:cka}) and the
  necessity of fan-in two (\cref{prop:fanin}); universality of label forests
  (\cref{thm:label}); merging trees gives mixtures (\cref{prop:interweave}).
\item One learned tree (\cref{sec:learn}): no operator deletes a leaf
  (\cref{lem:leaves}), so a learned tree is universally consistent whatever structure
  category utility builds (\cref{thm:consistency}).
\item Cognitive tasks (\cref{sec:tasks}): tasks as cue--outcome pairs with observable
  outcomes (\cref{def:task}); a learned forest converges to the ideal forest of its tasks
  (\cref{thm:ideal}) and is consistent if and only if the tasks are sufficient
  (\cref{cor:iff}); funnel tasks are sufficient for monotone-additive targets
  (\cref{prop:suffclass}) but not for $x_1x_2$ (\cref{prop:xy}); cross-view tasks are
  sufficient under multi-view redundancy (\cref{prop:multiview}); categorization is
  always sufficient, at a dimension-dependent rate (\cref{thm:labelcons}).
\item Inferring tasks (\cref{sec:infer}): no rule that looks at a tree's own data can
  infer its task (\cref{thm:locality}, \cref{cor:parity}); verifying candidate tasks at
  the composing tree selects a sufficient assignment with probability tending to one
  (\cref{thm:selection}); this yields task operators chosen like Cobweb's own
  (\cref{proc:induction}).
\end{itemize}

%=====================================================================
\section{The model}\label{sec:model}
%=====================================================================

\subsection{Concepts and trees}

An \emph{instance} is a pair $z=(a,u)$ of a nominal part $a$ (a table that assigns to
each attribute a bag of values with weights) and a continuous part $u\in\R^m$.

\begin{definition}[Concept]\label{def:concept}
A concept $c$ stores its count $n_c$, the attribute--value counts $n_c(A{=}v)$ of the
instances it covers, and, for the continuous block, their mean $\mu_c\in\R^m$ and
empirical (per-dimension) variance $s_c^2\in\R^m_{\ge0}$.  With smoothing parameters
$\alpha>0$ and $h^2>0$ (\texttt{alpha}, \texttt{prior\_var}) a concept scores a query
with
\[
  P_c(A{=}v)=\frac{n_c(A{=}v)+\alpha}{n_c(A)+(K_A+1)\alpha},\qquad
  p_c(u_j)=\N\bigl(u_j;\,\mu_{c,j},\,v_{c,j}\bigr),\quad v_{c,j}=s^2_{c,j}+h^2,
\]
with attributes conditionally independent given $c$ and $K_A$ the number of values of $A$
seen so far (the extra $\alpha$ reserves mass for an unseen value).  When a concept
\emph{predicts} a missing attribute it uses $(n_c(A{=}v)+\alpha)/(n_c(A)+K_A\alpha)$,
which sums to one over the seen values.
\end{definition}

\begin{definition}[Cobweb tree]
A Cobweb tree $\T$ is a rooted tree of concepts whose statistics are \emph{additive}:
each internal concept's statistics are those of the pooled instances of its children.
Its \emph{category utility} at a node with children $C$ is the normalised information
gain $\mathrm{PU}=|C|^{-1}\bigl[H(\text{parent})-\sum_{c\in C}P(c)H(c)\bigr]$, where
$H$ is the sum of the Dirichlet-smoothed categorical entropies and of the Gaussian
differential entropies $\tfrac12\log(2\pi e\,v_{c,j})$; without the $|C|^{-1}$ factor
this is the mutual information $I(Z;C)$ \citep{corter1992,maclellan2026rational}.
\end{definition}

Learning (\texttt{ifit}) sorts an instance from the root.  At a leaf it either
increments the leaf (exact match) or performs a \emph{fringe split}: the leaf becomes an
internal node with the old leaf and a new leaf for the instance as children.  At an
internal node it adds the instance and chooses, by category utility, between
\textsc{best} (descend into the best child), \textsc{new} (create a leaf),
\textsc{merge} (join the two best children under a new node) and \textsc{split}
(replace the best child by its children).

\subsection{Prediction}

A \emph{query} is an instance in which some attributes and some continuous
coordinates are missing; we write $O$ for the observed coordinates (missing continuous
coordinates are encoded as \texttt{NaN}, missing attributes are absent).  The
log-likelihood of a query under a concept is
$\ell_c(q)=\sum_{A\text{ obs.}}\log P_c(A{=}q_A)+\sum_{j\in O}\log p_c(q_j)$.

\begin{definition}[Multi-node prediction]\label{def:predict}
Prediction expands concepts best-first from the root, scoring each non-root concept $c$
with one of
\[
\begin{array}{ll}
\textsc{lik}: & s_c=\ell_c(q),\\
\textsc{prior}: & s_c=\ell_c(q)+\log(n_c/n_{\text{root}}),\\
\textsc{path}: & s_c=\pi_c(q)+\ell_c(q),\qquad
  \pi_c(q)=\sum_{c'\in\text{path}(c)\setminus\{\text{root}\}}
  \log\frac{n_{c'}e^{\ell_{c'}(q)}}{\sum_{b\in\text{sib}(c')}n_b e^{\ell_b(q)}},\\
\textsc{posterior}: & s_c=\pi_c(q),
\end{array}
\]
where $\text{sib}(c')$ are the children of $c'$'s parent (including $c'$); under every
weighting the root is scored by its likelihood, $s_{\text{root}}=\ell_{\text{root}}(q)$.
With $E_M(q)$ the expanded set and $w_c(q)\propto e^{s_c}$ normalised over $E_M(q)$, the
prediction of a missing continuous coordinate $j$ and of a missing attribute $A$ is
\[
  \hy_j(q)=\sum_{c\in E_M(q)}w_c(q)\,\mu_{c,j},\qquad
  \hat P(A{=}v\mid q)=\sum_{c\in E_M(q)}w_c(q)\,P_c(A{=}v).
\]
The budget $M$ bounds the number of expanded concepts.  The \emph{leaves-only} variant
restricts the mixture to the expanded leaves; its budget counts leaves (expansion stops
once $M$ leaves have entered the mixture, after at most $64M$ expansions).  \emph{Full
expansion} is $M=\infty$.
\end{definition}

\textsc{path} is the default of the implementation and of \texttt{cobweb\_torch};
$\pi_c(q)$ is the log posterior of $c$ obtained by multiplying the sibling-normalised
posteriors down its path, so $\sum_{c\text{ at depth }k}e^{\pi_c}\le1$ for every
depth~$k$.  The rational account \citep{maclellan2026rational} ranks concepts by
pointwise mutual information $\log p(x\mid c)-\log p(x)$, i.e.\ by $\ell_c$.

\subsection{Cobweb forests}

\begin{definition}[Messages]\label{def:messages}
Let $\T$ be a tree whose instances contain \emph{outcome} attributes $V_\T$ (possibly
none; \cref{def:task}).  For a query $q$ in which $V_\T$ is missing:
\begin{itemize}[leftmargin=*]
\item the \emph{activation} of $\T$ is its prediction of $V_\T$---the vector
  $\hy_{V_\T}(q)$ for a continuous outcome, the predicted distribution
  $(\hat P(V_\T{=}v\mid q))_v$ for a nominal one;
\item the \emph{label} of $\T$ at depth $k$ is the identifier of the depth-$k$ concept
  on the root-to-leaf descent that chooses, at each node, the child maximising
  $\log n_c+\ell_c(q)$ (the leaf if the path is shorter; $k=\infty$ denotes the leaf).
\end{itemize}
\end{definition}

\begin{definition}[Cobweb forest]\label{def:forest}
A Cobweb forest $\F=(V,E,o)$ is a finite directed acyclic graph whose vertices are
Cobweb trees and whose edges $\sigma\to\tau$ carry the activation and/or the label of
$\sigma$ into $\tau$, together with an output tree $o\in V$.  A tree $\tau$ learns
instances made of (a)~the raw coordinates $x_{S_\tau}$ it is wired to, (b)~the activations
of its parents as continuous coordinates, (c)~the labels of its parents as nominal
attributes (one attribute per parent), and (d)~its outcome, if any.  For an input $x$
the forest computes the messages of every tree in topological order and outputs the
activation of $o$.  The \emph{fan-in} of $\tau$ is the number of its non-outcome inputs
(raw coordinates, activation coordinates, label attributes).
\end{definition}

\begin{definition}[Training protocols]\label{def:forestlearn}
Every tree learns only with \texttt{ifit}.  Two protocols fix the messages a tree is
trained on:
\begin{itemize}[leftmargin=*]
\item \emph{online}: for each observation $i$, every tree first emits its message for
  $x_i$ from what it has learned from observations $1,\dots,i-1$, then learns
  observation $i$ (messages are out-of-sample, as at test time);
\item \emph{split}: the stream is cut into consecutive chunks, one per level of the DAG;
  a tree at level $k$ learns chunk $k$, reading the messages of its already-trained
  (frozen) parents.
\end{itemize}
\end{definition}

\begin{remark}[Why not train stage-wise on the same data]
A tree that has learned $(x_i,y_i)$ keeps a leaf for it (\cref{lem:leaves}), and its
activation at $x_i$ is close to $y_i$ itself.  A parent trained on such in-sample
messages learns the relation between $y$ and a leaked copy of $y$, which it never sees at
test time.  The online protocol (predict, then learn), already used by the Cobweb
language model, and the split protocol avoid this.
\end{remark}

%=====================================================================
\section{Representation}\label{sec:repr}
%=====================================================================

Throughout, a tree that \emph{represents} $g:K\to\R$ ($K\subset\R^d$ compact) is a tree
over instances $(x,g(x))$ queried with the last coordinate missing, and
$\omega_g(r)=\sup\{|g(x)-g(x')|:\|x-x'\|\le r\}$ is the modulus of continuity.

\subsection{A local error certificate}

\begin{lemma}[Concept bias]\label{lem:bias}
Let $c$ be a concept of a tree over instances $(x_i,g(x_i))$ with $g$ $L$-Lipschitz.
Then $|\mu_{c,Y}-g(\mu_{c,X})|\le L\sqrt{\tr S_c}$, where $S_c$ is the empirical
covariance of the $x$-part of $c$'s instances, and consequently, for every query $x$,
every weighting and every budget,
\begin{equation}\label{eq:cert}
  |\hy(x)-g(x)|\;\le\;L\sum_{c\in E_M(x)}w_c(x)\Bigl(\|x-\mu_{c,X}\|+\sqrt{\tr S_c}\Bigr).
\end{equation}
\end{lemma}
\begin{proof}
$|\mu_{c,Y}-g(\mu_{c,X})|=|\tfrac{1}{n_c}\sum_i(g(x_i)-g(\mu_{c,X}))|\le
\tfrac{L}{n_c}\sum_i\|x_i-\mu_{c,X}\|\le L(\tfrac1{n_c}\sum_i\|x_i-\mu_{c,X}\|^2)^{1/2}$
by Jensen.  For~\eqref{eq:cert}, write $\hy-g(x)=\sum_c w_c(\mu_{c,Y}-g(\mu_{c,X}))+
\sum_c w_c(g(\mu_{c,X})-g(x))$ and use $\sum_cw_c=1$.
\end{proof}

The right side of \eqref{eq:cert} is computable from the tree: a mixture is accurate
when it puts its weight on concepts that are near the query \emph{and} tight.

\subsection{One tree}

\begin{theorem}[A single tree is a universal approximator]\label{thm:single}
Let $g\in C(K)$, $K\subset\R^d$ compact, and $\varepsilon>0$.  There is a Cobweb tree $\T$
and a variance floor $h^2>0$ such that $\sup_{x\in K}|\hy_\T(x)-g(x)|<\varepsilon$ under
full expansion with each of the weightings \textsc{lik} and \textsc{path}, and with the
leaves-only variant of any weighting.  If $g$ is $L$-Lipschitz, a tree with $G$ leaves
achieves
\begin{equation}\label{eq:singlerate}
  \sup_{x\in K}|\hy_\T(x)-g(x)|\;\le\;C_d\,L\,G^{-1/d}\sqrt{\log G}+C_d\|g\|_\infty G^{-1}.
\end{equation}
\end{theorem}

\begin{proof}
Extend $g$ continuously to a cube $Q\supseteq K$ (Tietze) and rescale $Q=[0,1]^d$.
Fix $\beta>0$, let $\delta=\beta h$ and let $\Gamma$ be the grid of cell centres of
spacing $\delta$, $G=|\Gamma|=\delta^{-d}$.  Take the depth-one tree whose leaves are
the instances $(\gamma,g(\gamma))$, $\gamma\in\Gamma$, and whose root pools them.  Each
leaf has $s^2=0$, hence variance $h^2$ in every coordinate; the root has variance at
least $v_0>0$ (the variance of the grid) in every coordinate.  Write
$k_\gamma=(2\pi h^2)^{-d/2}e^{-\|x-\gamma\|^2/2h^2}$.

\emph{Leaves dominate.}  Some $\gamma$ lies within $\sqrt d\,\delta/2$ of $x$, so
$\max_\gamma k_\gamma\ge(2\pi h^2)^{-d/2}e^{-d\beta^2/8}$, while
$\sum_\gamma k_\gamma\le(2\pi h^2)^{-d/2}C(\beta,d)$ by comparison with a Gaussian sum
over a grid of spacing $\beta$ in units of $h$.  Under \textsc{lik} the leaves carry
unnormalised mass $\sum_\gamma k_\gamma\ge c\,h^{-d}$ and the root at most
$(2\pi v_0)^{-d/2}$.  Under \textsc{path} the leaf $\gamma$ has
$e^{\pi_\gamma}=k_\gamma/\sum_{\gamma'}k_{\gamma'}$ (equal counts), so its mass is
$k_\gamma^2/\sum k$, whose sum is at least $(\max k)^2/\sum k\ge c'h^{-d}$, while the root
again has mass at most $(2\pi v_0)^{-d/2}$.  In both cases the root's weight is
$O(h^d)=O(G^{-1})$.

\emph{Leaves form a kernel smoother.}  Among the leaves, $w_\gamma\propto k_\gamma^{\,r}$
with $r=1$ (\textsc{lik}) or $r=2$ (\textsc{path}).  For any $\rho>0$, splitting the
exponent and using the lower bound on $\max_\gamma k_\gamma^r$ for the normaliser,
\[
\sum_{\|\gamma-x\|>\rho}w_\gamma\;\le\;
\frac{e^{-r\rho^2/4h^2}\sum_\gamma e^{-r\|x-\gamma\|^2/4h^2}}{e^{-rd\beta^2/8}}
\;\le\;C''e^{-r\rho^2/4h^2},
\]
so $|\sum_\gamma w_\gamma g(\gamma)-g(x)|\le\omega_g(\rho)+2\|g\|_\infty C''e^{-\rho^2/4h^2}$.
Taking $\rho=\sqrt h$ makes this vanish as $h\to0$, uniformly in $x$; adding the root's
$O(h^d)$ weight (times $2\|g\|_\infty$) gives the first claim.  For Lipschitz $g$ take
$\rho=h\sqrt{2\log(1/h)}$, so $\omega_g(\rho)\le L\rho$ and the tail is $O(h)$; with
$h\asymp G^{-1/d}$ this is \eqref{eq:singlerate}.  The leaves-only
variant drops the root term; restricted to the leaves of this tree, \textsc{prior}
coincides with \textsc{lik} (equal counts) and \textsc{posterior} weights the leaves by
$k_\gamma/\sum k$, i.e.\ $r=1$.
\end{proof}

The \textsc{prior} weighting behaves differently: it multiplies each concept's
likelihood by its \emph{relative count}, and the internal concepts then keep a share of
the mixture that does not vanish as the variance floor shrinks.

\begin{proposition}[The prior- and posterior-weighted all-concept mixtures are inconsistent]\label{prop:prior}
In the setting of the proof of \cref{thm:single}, with \textsc{prior} weighting and full
expansion, the root's weight satisfies
$\lambda(x)\to\N_{\text{root}}(x)/(1+\N_{\text{root}}(x))$ as $h\to0$ for $x$ in the interior
of $[0,1]^d$, where $\N_{\text{root}}(x)=\N(x;\mu_{\text{root}},v_{\text{root}})$, and $\hy(x)-g(x)=\lambda(x)\bigl(\mu_{\text{root},Y}-g(x)\bigr)+o(1)$: the
error does not vanish wherever $g(x)$ differs from the grid mean of $g$.  With
\textsc{posterior} weighting the leaves' masses $e^{\pi_\gamma}$ sum to exactly one, so
$\lambda(x)=\N_{\text{root}}(x)/(\N_{\text{root}}(x)+1)$ and the same conclusion holds.
\end{proposition}
\begin{proof}
Leaf $\gamma$ has prior $1/G$, so the leaves' mass is $G^{-1}\sum_\gamma k_\gamma$, a
Riemann sum of the Gaussian kernel density estimate of the uniform density on $[0,1]^d$,
which tends to $1$ for $x$ in the interior as $h\to0$.  The root has prior $1$ and mass $\N_{\text{root}}(x)$, which tends to a positive limit.
The leaves' mixture converges to $g(x)$ as in \cref{thm:single}.  Under \textsc{posterior}
the leaves are the root's children and their path posteriors are the sibling-normalised
posteriors, which sum to one, while the root is scored by its likelihood
(\cref{def:predict}).
\end{proof}

In a learned tree the defect of \textsc{posterior} is compounded: every non-root concept
is scored by its path posterior $\pi_c\le0$, which can only decrease down a path, so the
wide concepts near the root outweigh the leaves whatever $h$ is.  Both are
\emph{all-concept} pathologies; restricted to leaves the weightings converge in the
construction above, and \cref{sec:experiments} measures them in learned trees.

\subsection{Composition: the Cobweb--Kolmogorov--Arnold theorem}

\begin{lemma}[Approximating a DAG of continuous maps]\label{lem:dag}
Let $D$ be a finite DAG of continuous maps $g_v:\R^{k_v}\to\R$ whose inputs are raw
coordinates of $x\in[0,1]^n$ or outputs of other vertices, and let $F$ be the map it
computes; let $K_v$ be the (compact) set of input vectors of $v$ as $x$ ranges over
$[0,1]^n$.  For every $\varepsilon>0$ there are $\rho_v,\eta_v>0$ such that any maps
$\hat g_v$ with $\sup_{K_v+\bar B_{\rho_v}}|\hat g_v-g_v|\le\eta_v$ define a network
$\hat F$ with $\sup_x|\hat F(x)-F(x)|\le\varepsilon$.
\end{lemma}
\begin{proof}
Choose the tolerances in reverse topological order.  At the output vertex put
$\eta_o=\varepsilon/2$ and choose $\rho_o\le1$ with
$\omega_{g_o}(\rho_o)\le\varepsilon/2$ on the compact $K_o+\bar B_1$.  Given the
tolerance $\rho_v\le1$ by which the input vector of $v$ may deviate, allow each of its
$k_v$ inputs to deviate by $\rho_v/\sqrt{k_v}$: for each incoming vertex $u$ set
$\eta_u=\rho_v/(2\sqrt{k_v})$ and choose $\rho_u\le1$ with
$\omega_{g_u}(\rho_u)\le\rho_v/(2\sqrt{k_v})$ on $K_u+\bar B_1$, taking minima when $u$
feeds several vertices.  Then, in topological order, every vertex's input deviates from
its exact value by at most $\rho_v$, so its output deviates by at most
$\eta_v+\omega_{g_v}(\rho_v)$, which is within its consumers' tolerance; at the output
the deviation is at most $\varepsilon$.
\end{proof}

\begin{theorem}[Cobweb--Kolmogorov--Arnold]\label{thm:cka}
Let $n\ge2$.  For every $f\in C([0,1]^n)$ and $\varepsilon>0$ there is a Cobweb forest
with activation messages, in which every tree has at most two inputs, consisting of
\begin{itemize}[leftmargin=*,noitemsep]
\item $n(2n+1)$ univariate trees for the inner functions $\varphi_{q,p}$ (independent of
  $f$),
\item $(2n+1)(n-1)$ two-input trees forming the inner sums,
\item $2n+1$ univariate trees for the outer functions $\Phi_q$, and
\item $2n$ two-input trees forming the outer sum,
\end{itemize}
$4n^2+4n$ trees in all, whose output is within $\varepsilon$ of $f$ on $[0,1]^n$.
\end{theorem}
\begin{proof}
Write $f$ in the form \eqref{eq:ka} and read it as a DAG of continuous maps of one or
two variables: $\varphi_{q,p}$, the binary additions $(a,b)\mapsto a+b$ chained to form
the $2n+1$ inner sums of $n$ terms and the outer sum of $2n+1$ terms, and $\Phi_q$.
Apply \cref{lem:dag}, and realise each $\hat g_v$ on the compact
$K_v+\bar B_{\rho_v}$ by \cref{thm:single} as a tree over instances
$(\text{inputs},g_v(\text{inputs}))$; its activation is the required $\hat g_v$.  The
count is $n(2n+1)+(2n+1)(n-1)+(2n+1)+2n=4n^2+4n$.
\end{proof}

\begin{proposition}[Two inputs are necessary]\label{prop:fanin}
If every tree of a forest has fan-in at most one, its output depends on at most one raw
coordinate.  Hence fan-in two is the least fan-in at which Cobweb forests are universal.
\end{proposition}
\begin{proof}
By induction in topological order, each tree's messages are functions of the at most one
raw coordinate at the source of its unique input chain.
\end{proof}

\begin{remark}[Rates and smooth KA representations]\label{rem:rates}
\Cref{thm:cka} uses a fixed number of trees, but for an arbitrary $f$ the inner functions
of \eqref{eq:ka} are rough, and the number of leaves each tree needs is not controlled.
As for KANs \citep[Thm.~2.1]{liu2024kan}, the useful statement assumes structure.  If
$f$ is computed by a DAG of $s$ maps of fan-in $\le k$, each $L$-Lipschitz on a
neighbourhood of its input set, then \eqref{eq:singlerate} and the stability argument of
\cref{lem:dag} give a forest of $s$ trees with $G$ leaves each and error
$C(s,L)\,G^{-1/k}\sqrt{\log G}$---a rate set by the fan-in $k$, not by $n$.  A single
tree over the $n$ inputs has rate $G^{-1/n}$ (\cref{thm:single}).  Cobweb trees realise
sums as learned two-input maps; KANs sum at nodes for free, which is why their rate is
governed by one-dimensional splines.
\end{remark}

\subsection{Label forests}

\begin{theorem}[Label forests are universal]\label{thm:label}
For every $f\in C([0,1]^n)$ and $\varepsilon>0$ there is a forest of $n$ univariate trees
without outcomes that pass leaf labels to an output tree whose instances consist of the
$n$ labels (nominal attributes) and the target, such that the output tree's prediction
(\textsc{lik}, leaves only, full expansion) is uniformly within $\varepsilon$ of $f$.
\end{theorem}
\begin{proof}
Let tree $p$ be the depth-one tree over the $K$ grid points of spacing $\delta=1/K$ in
$[0,1]$, all with count one and variance $h^2$.  Its label for $x_p$ is
$\arg\max_\gamma\{\log n_\gamma+\ell_\gamma(x_p)\}$, the nearest grid point, so the
labels partition $[0,1]^n$ into boxes $B_t$ of side $\delta$, indexed by label tuples
$t$.  Let the output tree be the depth-one tree with one leaf per tuple $t$, holding the
instance $(t,f(\text{centre}(B_t)))$.  For a query with tuple $t$, a leaf $t'$ that
disagrees with $t$ in $m$ attributes has likelihood ratio $\bigl(\alpha/(1+\alpha)\bigr)^m$
relative to leaf $t$ (\cref{def:concept}, counts one).  Since each such leaf has weight at most $\alpha^m$ relative to leaf $t$, their total
weight is at most $\sum_{m\ge1}\binom nm\alpha^m=(1+\alpha)^n-1$.  Hence
$|\hy(x)-f(x)|\le\omega_f(\sqrt n\,\delta/2)+2\|f\|_\infty\bigl((1+\alpha)^n-1\bigr)$,
which is below $\varepsilon$ for $K$ large and then $\alpha$ small.
\end{proof}

The label forest of \cref{thm:label} has an output tree with $K^n$ leaves: labels carry
no metric, so this construction pays for every cell of the product partition.
Activations carry the metric and allow the fan-in-two construction of \cref{thm:cka}.

\subsection{Interweaving trees}

\begin{proposition}[Merging trees under a common root gives mixtures]\label{prop:interweave}
Let $\T_1,\dots,\T_m$ be trees over the same instance space with the same value sets
(e.g.\ continuous-only trees), and let $\T$ be the tree whose root pools their roots'
statistics and has $\T_1,\dots,\T_m$ as its subtrees.  Under full expansion with
\textsc{lik} or \textsc{path} weighting,
\[
  \hy_\T(q)=\lambda_0(q)\,\mu_{\text{root}}+\sum_{j=1}^m\lambda_j(q)\,\hy_{\T_j}(q),
  \qquad\lambda_j\ge0,\ \sum_{j=0}^m\lambda_j=1,
\]
with $\lambda_0\propto e^{\ell_{\text{root}}}$ and $\lambda_j\propto Z_j$ (\textsc{lik})
or $\lambda_j\propto P(\text{root}_j\mid q)\,Z_j$ (\textsc{path}), where $Z_j$ is the
total unnormalised weight of $\T_j$'s own mixture and $P(\text{root}_j\mid q)$ the
sibling-normalised posterior of $\T_j$'s root.
\end{proposition}
\begin{proof}
Under \textsc{lik} a concept's score does not depend on the tree it sits in, so the
mixture over $\T$ splits into the root term and the $m$ subtree mixtures.  Under
\textsc{path} every concept of $\T_j$ gains the same term $\log P(\text{root}_j\mid q)$
in $\pi_c$ (the sibling sets below $\text{root}_j$ are unchanged), which rescales the
subtree's weights uniformly and leaves its internal mixture unchanged.  With different
value sets the smoothing denominators $K_A$ change, and the identity holds only
approximately.
\end{proof}

Interweaving therefore closes the class of tree predictors under input-dependent
convex combinations; it never composes them.  Composition---one tree reading another's
representation---is what forests add.

%=====================================================================
\section{Learnability of one tree}\label{sec:learn}
%=====================================================================

\begin{lemma}[Leaf preservation]\label{lem:leaves}
Run \texttt{ifit} on any sequence $z_1,\dots,z_N$ with any choices of operators
(including add-only training), and suppose the depth limit is never reached.  Then
(i)~every concept's statistics are the pooled statistics of the instances stored in the
leaves below it; (ii)~no leaf is ever deleted; (iii)~the leaves are in bijection with the distinct
instances, where two instances are the same if they agree up to the exact-match
tolerance.
\end{lemma}
\begin{proof}
By induction over the branches of the sorting loop.  \textsc{best} increments the
current node and descends; \textsc{new} increments the current node and creates a leaf
holding $z$; \textsc{merge} inserts a new internal node whose statistics are the pooled
statistics of the two children it adopts; \textsc{split} deletes the best child, which is
internal (the operator requires it to have at least two children), and promotes its
children; at a leaf, an exact match increments it and any other instance triggers a
fringe split, which copies the leaf's statistics into a new internal parent, keeps the
old leaf, and adds a leaf for $z$.  No branch deletes a leaf, every instance ends in
exactly one leaf, and a leaf only absorbs exact matches.
\end{proof}

When the depth limit is reached the implementation inserts the instance into the
current node, which may be a leaf; (iii) then fails, which is why the lemma assumes the
default limit (1000) is not reached.

\begin{theorem}[A learned tree is universally consistent]\label{thm:consistency}
Let $(X_i,Y_i)$ be i.i.d.\ with $X\in\R^d$, $|Y|\le B$, $m(x)=\E[Y\mid X=x]$.
Let $\T_N$ be the tree that Cobweb learns from $(X_i,Y_i)_{i\le N}$ in any order and with
any operator choices, and let it predict $Y$ from $X$ with the leaves-only mixture,
\textsc{lik} weighting and full expansion, with variance floor $h_N^2$.  If $X$ has a
density, $h_N\to0$ and $Nh_N^d\to\infty$, then
$\int|\hy_{\T_N}(x)-m(x)|\,\mu(dx)\to0$ almost surely, whatever structure category
utility builds ($\mu$ the distribution of $X$).
\end{theorem}
\begin{proof}
With a density the $X_i$ are distinct almost surely and so are the instances, so by
\cref{lem:leaves} the tree has exactly $N$ leaves, leaf $i$ holding $(X_i,Y_i)$ with zero
empirical variance.  The \textsc{lik} score of leaf $i$ for a query $x$ with $Y$
missing is $\log\N(x;X_i,h_N^2I)$, so the leaves-only mixture is
$\sum_iK_{h_N}(x-X_i)Y_i/\sum_iK_{h_N}(x-X_i)$, the Nadaraya--Watson estimate with
a Gaussian kernel.  The Gaussian kernel is \emph{regular} (it is bounded below by a
positive constant on a ball and $\int\sup_{\|u-x\|\le r}K(u)\,dx<\infty$), and for
regular kernels and bounded $Y$ this estimate is strongly consistent in $L_1(\mu)$
under $h_N\to0$, $Nh_N^d\to\infty$ \citep{devroye1989}; see \citet[Ch.~5]{gyorfi2002}
for the boxed-kernel version in $L_2$.
\end{proof}

\begin{remark}[What the structure is for]\label{rem:structure}
\Cref{thm:consistency} holds because Cobweb never forgets an instance; the hierarchy
plays no role in it.  The structure matters for three things the theorem does not
cover: (a)~a \emph{budget} $M<\infty$, where best-first search must find the leaves
near the query; (b)~the \textsc{path} weighting, where the mixture weight of a leaf is
its path posterior times its likelihood, so the leaves-only mixture is not a kernel
smoother in a learned tree (\cref{lem:bias} still bounds its error); and (c)~the
\emph{labels} a tree passes to other trees, which are its concepts.  The experiments use
the implementation's default, \textsc{path} with a leaf budget; they are therefore a
test of these structural effects as much as of the theorem.  The branch-and-bound stop
of \texttt{categorize(use\_best=True)} assumes $\ell_c\le0$, which holds for nominal
attributes but not for Gaussian densities; with a continuous block it can return a
concept whose score is below the best (\cref{app:verify} exhibits one).
\end{remark}

%=====================================================================
\section{Cognitive tasks}\label{sec:tasks}
%=====================================================================

A Cobweb tree is unsupervised: it models the joint distribution of whatever its
instances contain and predicts any attribute that is missing from a query.  What a tree
\emph{passes on} is therefore fixed by what its instances contain.  We make this the
central object.

\subsection{Tasks as cues and outcomes}

\begin{definition}[Cognitive task]\label{def:task}
Let the stream consist of i.i.d.\ observations $(X,Y)$, $X\in\R^n$ the raw input and
$Y$ the target.  A \emph{cognitive task} for a tree $\tau$ is a pair $(C_\tau,V_\tau)$:
\begin{itemize}[leftmargin=*]
\item the \emph{cue} $C_\tau$, what the tree perceives: its raw coordinates
  $X_{S_\tau}$ and the messages of its parents;
\item the \emph{outcome} $V_\tau$, what the tree predicts and later observes: either
  nothing ($V_\tau=\varnothing$) or an \emph{observable} variable, i.e.\ a coordinate of
  the observation that the cue does not contain---the target $Y$ or a raw coordinate
  $X_j$, $j\notin S_\tau$.
\end{itemize}
The tree learns the instances $(C_\tau,V_\tau)$ with \texttt{ifit} and is queried with
$V_\tau$ missing (\cref{def:messages}).  A \emph{task assignment} $A$ gives every tree
of a forest a task; the output tree's outcome is $Y$.
\end{definition}

Three families of tasks occur in the systems that motivate this paper.
\emph{Categorization} ($V_\tau=\varnothing$): the tree passes only labels, as the filter
hierarchy of Convolutional Cobweb does.  \emph{Funnel} ($V_\tau=Y$): every tree predicts
the final target, as every tree of the Cobweb language model predicts the next word.
\emph{Cross-view} ($V_\tau=X_j$, $j\notin S_\tau$): the tree predicts a part of the input
it does not perceive---a neighbouring pixel, the next token, the context of a chunk, as
TRELLIS's content hierarchy models its context.  The cue--outcome reading is the one of
associative learning, in which an organism learns to predict an outcome it will observe
from the cues it perceives \citep{rescorla1972}; the online protocol of
\cref{def:forestlearn} is literally ``predict, then observe, then learn''.

Observability is the constraint that separates learnable from representable forests.

\begin{remark}[The KA tasks are not observable]\label{rem:kaobs}
In \cref{thm:cka} the tree for $\Phi_q$ holds the values $\Phi_q(u_q)$.  The stream
never shows them: it shows only their sum $Y=\sum_q\Phi_q(u_q)$.  A KAN separates the
terms of the sum by back-propagating the error of the sum through each of them; a Cobweb
tree can only predict what it will observe.  The inner functions $\varphi_{q,p}$ do not
depend on $f$, so a forest could be \emph{given} them as fixed tasks, but the outer
decomposition is not identifiable from the stream without further structure.  The rest
of this section asks which \emph{observable} tasks recover $f$ anyway.
\end{remark}

\subsection{What a learned forest converges to}

\begin{definition}[Ideal forest, sufficiency, deficit]\label{def:ideal}
Let $A$ be a task assignment of an activation forest.  Define, in topological order,
the \emph{ideal cue} and \emph{ideal message} of each tree,
\[
  U_\tau=\bigl(X_{S_\tau},(m_\sigma(X))_{\sigma\to\tau}\bigr),\qquad
  m_\tau(X)=\E[V_\tau\mid U_\tau]=\psi_\tau(U_\tau),
\]
and the \emph{ideal output} $F_A(X)=m_o(X)=\E[Y\mid U_o]$.  Let $f(X)=\E[Y\mid X]$.  The
assignment is \emph{sufficient} if $F_A=f$ almost surely, and its \emph{deficit} is
\[
  \Delta(A)=\E\bigl[(f(X)-F_A(X))^2\bigr]=\E\Var(Y\mid U_o)-\E\Var(Y\mid X).
\]
\end{definition}

The identity holds because $U_o$ is a function of $X$, so $F_A=\E[f(X)\mid U_o]$ and
$\Var(Y\mid U_o)=\E[\Var(Y\mid X)\mid U_o]+\Var(f(X)\mid U_o)$.

\begin{assumption}[Regularity]\label{ass:regular}
(A1) The observations are i.i.d., $X$ has compact support, every outcome is bounded,
$|V_\tau|\le B$, and the instances of every tree are almost surely pairwise distinct
(e.g.\ $V_\tau$ has a continuous distribution).  (A2) For every tree, the ideal cue
$U_\tau\in\R^{k_\tau}$ has a density bounded away from zero and infinity on its support
$D_\tau$, a compact set with Lipschitz boundary, and $\psi_\tau$ extends to a uniformly
continuous function on a neighbourhood of $D_\tau$.  (A3) Every tree predicts with the leaves-only \textsc{lik}
mixture under full expansion, with variance floor $h_{\tau,N}^2$, $h_{\tau,N}\to0$ and
$Nh_{\tau,N}^{k_\tau}/\log N\to\infty$.  (A4) For every edge $\sigma\to\tau$,
$\varepsilon_{\sigma,N}/h_{\tau,N}\to0$ in probability, where
$\varepsilon_{\sigma,N}=\sup_x|\hat m_\sigma(x)-m_\sigma(x)|$ is the parent's uniform
error.
\end{assumption}

(A4) is a condition on the schedule: by induction the parents' errors tend to zero
(\cref{thm:ideal}), so variance floors that shrink slowly enough satisfy it.

\begin{theorem}[A Cobweb-learned forest converges to its ideal forest]\label{thm:ideal}
Let $\F$ be an activation forest with task assignment $A$, trained with the split
protocol on $N$ observations, and let \cref{ass:regular} hold.  Then for every tree
$\tau$, $\sup_x|\hat m_\tau(x)-m_\tau(x)|\to0$ in probability; in particular the output
satisfies $\sup_x|\hat F_N(x)-F_A(x)|\to0$ in probability.
\end{theorem}
\begin{proof}
Induction over the levels of the DAG.  Fix a tree $\tau$ at level $\ell$ and condition on
the chunks of the earlier levels, on which its parents were trained and frozen.  Write
$\hat u_\tau(x)=(x_{S_\tau},(\hat m_\sigma(x))_\sigma)$ for its learned cue map and
$u_\tau(x)$ for the ideal one, so $\|\hat u_\tau-u_\tau\|_\infty\le\varepsilon:=
\sqrt{k_\tau}\max_\sigma\varepsilon_{\sigma,N}$.

\emph{Step 1: the tree is a kernel smoother.}  The tree learns the instances
$(\hat U_i,V_i)=(\hat u_\tau(X_i),V_{\tau,i})$ of chunk $\ell$, which are i.i.d.\ given
the earlier chunks and pairwise distinct by (A1).  By \cref{lem:leaves} its leaves are
these instances with zero spread, and its leaves-only \textsc{lik} activation at a cue
$u$ is the Nadaraya--Watson estimate
$\hat\psi(u)=\sum_iK_h(u-\hat U_i)V_i/\sum_iK_h(u-\hat U_i)$ with the Gaussian kernel of
bandwidth $h=h_{\tau,N}$.

\emph{Step 2: the perturbed design does not matter.}  Let $\check\psi$ be the same
estimate built on the ideal cues $U_i=u_\tau(X_i)$.  By (A1)--(A3), $\check\psi\to\psi_\tau$
uniformly on the $\delta$-neighbourhood of $D_\tau$ for any $\delta=o(h)$, almost surely
\citep{devroye1978,mack1982}.  Write $\hat U_i=U_i+\Delta_i$, $\|\Delta_i\|\le\varepsilon$.
For $\|u-U_i\|\le Rh$,
\[
  \frac{K_h(u-\hat U_i)}{K_h(u-U_i)}=\exp\Bigl(\frac{2\langle u-U_i,\Delta_i\rangle-\|\Delta_i\|^2}{2h^2}\Bigr)
  \in\bigl[e^{-\eta},e^{\eta}\bigr],\qquad \eta=\frac{R\varepsilon}{h}+\frac{\varepsilon^2}{2h^2},
\]
and for $\|u-U_i\|>Rh\ge2\varepsilon$, $K_h(u-\hat U_i)\le(2\pi h^2)^{-k/2}
e^{-\|u-U_i\|^2/8h^2}$.  Since the density of $U_\tau$ is bounded below on $D_\tau$ and
bounded above, and $D_\tau$ has Lipschitz boundary, the kernel mass of the near points
is at least $c\,N$ and that of the far points at most $C\,Ne^{-R^2/16}$, uniformly in $u$
within $o(h)$ of $D_\tau$, with probability tending to one (uniform laws of large numbers
for kernel classes).
Hence $\sup_u|\hat\psi(u)-\check\psi(u)|\le2B\bigl(e^{2\eta}-1+C'e^{-R^2/16}\bigr)$, which
tends to zero in probability when $R=R_N\to\infty$ slowly enough that
$R\varepsilon/h\to0$, possible by (A4).

\emph{Step 3: the test-time cue.}  At test time the tree is queried at $\hat u_\tau(x)$,
so $|\hat m_\tau(x)-m_\tau(x)|\le\sup_{D_\tau^{(\varepsilon)}}|\hat\psi-\psi_\tau|+
\omega_{\psi_\tau}(\varepsilon)$, where $D_\tau^{(\varepsilon)}$ is the
$\varepsilon$-neighbourhood of $D_\tau$; both terms tend to zero in probability by
Steps~1--2 and (A2).  This is the induction hypothesis for the children of $\tau$; at
the output tree it is the claim.
\end{proof}

\begin{corollary}[Consistency if and only if sufficiency]\label{cor:iff}
Under the hypotheses of \cref{thm:ideal}, $\E[(\hat F_N(X)-f(X))^2]\to\Delta(A)$.  The
forest is consistent for $f$ if and only if its task assignment is sufficient.
\end{corollary}
\begin{proof}
$\hat F_N$ and $f$ are bounded by $B$, so uniform convergence in probability gives
$\E(\hat F_N-F_A)^2\to0$, and $\E(\hat F_N-f)^2\to\E(F_A-f)^2=\Delta(A)$.
\end{proof}

\Cref{thm:ideal} separates what Cobweb's learning does from what the tasks determine:
every tree learns the conditional expectation of its outcome given its cue, nothing more
and nothing less; whether the composition of these conditional expectations is $f$ is a
property of the task assignment alone.

\subsection{Which tasks are sufficient}

\paragraph{Funnel tasks.}
A convenient sufficient condition for funnel forests is that each tree's ideal message
equals $\E[Y\mid X_{R(\tau)}]$, where $R(\tau)$ is the set of raw coordinates upstream of
$\tau$; this holds by induction when every $\E[Y\mid X_{R(\tau)}]$ is a function of the
parents' $\E[Y\mid X_{R(\sigma)}]$ and of $X_{S_\tau}$.

\begin{proposition}[Funnel tasks are sufficient for monotone-additive targets]\label{prop:suffclass}
Let $X$ consist of independent blocks $X_{B_1},\dots,X_{B_m}$, and let
$Y=\Phi\bigl(\sum_jg_j(X_{B_j})\bigr)+\xi$ with $\Phi$ continuous,
$g_j$ continuous, and $\xi$ centred and independent of $X$.  Then every funnel
activation forest whose graph is a rooted tree in which each block is the raw cue of
exactly one tree is sufficient.
\end{proposition}
\begin{proof}
For an upstream set of blocks $S$, $\E[Y\mid X_S]=\E[\Phi(s_S+Z_{\bar S})]=:H_S(s_S)$ with
$s_S=\sum_{j\in S}g_j(X_{B_j})$ and $Z_{\bar S}$ the sum over the other blocks, independent
of $X_S$; $H_S$ is continuous and injective because $\Phi$ is continuous.  A tree whose
children cover the disjoint block sets $S_1,\dots,S_k$ and which reads the raw blocks $S_0$
has $s_S=\sum_iH_{S_i}^{-1}(m_{S_i})+\sum_{j\in S_0}g_j(X_{B_j})$, so its
$\E[Y\mid X_S]=H_S(s_S)$ is a continuous function of its parents' messages and its raw
cue.
\end{proof}

The KAN examples $\exp(\sin\pi x_1+x_2^2)$, $\exp\bigl(\tfrac1{100}\sum_i\sin^2(\pi x_i/2)\bigr)$
and $\exp\bigl(\tfrac12(\sin\pi(x_1^2+x_2^2)+\sin\pi(x_3^2+x_4^2))\bigr)$ (the last with
the blocks $(x_1,x_2)$ and $(x_3,x_4)$) are in this class.

\begin{proposition}[Funnel tasks can fail]\label{prop:xy}
Let $X_1,X_2$ be independent and symmetric about $0$, and $Y=X_1X_2$.  For the funnel
forest whose output tree reads the activations of a tree perceiving $X_1$ and a tree
perceiving $X_2$, $m_1=m_2\equiv0$, $F_A\equiv0$ and $\Delta(A)=\Var(Y)$: the forest
learns nothing, whatever $N$ and whatever structure its trees build.
\end{proposition}
\begin{proof}
$\E[X_1X_2\mid X_1]=X_1\E X_2=0$, and likewise for $X_2$; the output tree's ideal cue is
constant.  By \cref{thm:consistency} the source trees' activations converge to $0$.
\end{proof}

\paragraph{Cross-view tasks.}
A tree that predicts another part of the input never sees the target.  It recovers what
the views share, which suffices when the target depends on the input only through what
the views share---the multi-view redundancy assumption of co-training and of
contrastive learning \citep{blum1998,tosh2021}.

\begin{proposition}[Cross-view tasks are sufficient under multi-view redundancy]\label{prop:multiview}
Let $Z\in\R^r$ be latent and $X=(X^{(1)},\dots,X^{(m)})$ consist of $m$ views with
$(Z,X)$ jointly Gaussian and centred, $X^{(v)}=\Lambda_vZ+E_v$, the noises $E_v$
independent of each other and of $Z$.  Let $Y=g(Z)+\xi$ with $\E|g(Z)|<\infty$ and $\xi$
centred and independent of $(Z,X)$.  Give the tree of view $v$ the cue $X^{(v)}$ and as
outcome $r$ coordinates $W_v$ of other views whose loading matrix $\Lambda_{W_v}\in\R^{r\times r}$
is invertible, and let the output tree have the activations of the $m$ view trees as cue
and $Y$ as outcome.  This assignment is sufficient, and no tree but the output tree
perceives $Y$.
\end{proposition}
\begin{proof}
Write $\mu_v=\E[Z\mid X^{(v)}]$ and $P_v$ for the posterior precision given $X^{(v)}$, and
$P_0$ for the prior precision; both precisions are constants.  Since $W_v=\Lambda_{W_v}Z+E'$
with $E'$ independent of $(Z,X^{(v)})$, the ideal message is
$m_v=\E[W_v\mid X^{(v)}]=\Lambda_{W_v}\mu_v$, an invertible linear function of $\mu_v$.
By conditional independence of the views,
$p(z\mid x)\propto p(z)\prod_vp(x^{(v)}\mid z)\propto p(z)^{1-m}\prod_vp(z\mid x^{(v)})$,
a Gaussian with precision $P=\sum_vP_v-mP_0$ and mean $P^{-1}\sum_vP_v\mu_v$.  The
posterior of $Z$ given $X$ therefore depends on $X$ only through $(\mu_v)_v$, hence through
$(m_v)_v$, and so does $f(X)=\E[g(Z)\mid X]$.  Thus $F_A=\E[Y\mid(m_v)_v]=f$.
\end{proof}

The assumption is necessary in the following sense: if the target carries a
view-specific component, e.g.\ $Y=e^\top E_1$ with $X^{(1)}$ of dimension larger than
$r$ and $e$ generic, the cross-view messages discard it and the deficit is positive.  In images and text the
views are patches and contexts, and the cross-view outcome is a neighbouring pixel or
the next token.

\paragraph{Categorization.}
A tree without an outcome passes labels.  Its labels do not depend on $f$ at all, which
is their strength and their cost.

\begin{theorem}[Consistency of label forests]\label{thm:labelcons}
Let a forest consist of trees without outcomes over the raw coordinates passing labels
at depth $k_N$ (or leaf labels) to an output tree that models $(\text{labels},Y)$, learned
with the split protocol, and let the output tree predict with the leaves-only
\textsc{lik} mixture with $\alpha_N\to0$ fast enough
($\alpha_N\cdot\#\text{labels per tree}\cdot n\to0$).  Let $\mathcal P_N$ be the product
partition of $[0,1]^n$ induced by the frozen label maps.  If (a)~the diameter of the cell
of $\mathcal P_N$ containing $X$ tends to $0$ in probability and (b)~the number of cells
of $\mathcal P_N$ that contain a sample of the output chunk is $o(N)$, then the forest is
a consistent regression estimate.
\end{theorem}
\begin{proof}[Proof sketch]
As in the proof of \cref{thm:label}, with $\alpha_N\to0$ the output tree's leaves-only
mixture for a query in cell $B$ puts weight $1-o(1)$ uniformly on the leaves whose
label tuple equals that of $B$ (one per training point of the output chunk in $B$, by
\cref{lem:leaves}); their $y$-means average to the partitioning estimate
$\sum_{X_i\in B}Y_i/\#\{X_i\in B\}$.  Conditional on the frozen partition this is the
classical partitioning estimate, consistent under (a) and (b) \citep[Thm.~4.2]{gyorfi2002}.
\end{proof}

\begin{remark}[The trade-off]\label{rem:tradeoff}
Condition (b) requires the product partition to have fewer occupied cells than samples;
with $K$ labels per tree that is $K^n=o(N)$, so the rate of a label forest is governed by
$n$, not by the fan-in.  Funnel and cross-view forests have rates governed by the fan-in
(\cref{thm:ideal} with \cref{rem:rates}) but need sufficiency (\cref{prop:xy}).
\end{remark}

\begin{conjecture}[Hybrid messages]\label{conj:hybrid}
A forest whose trees pass labels and funnel activations is consistent for every
continuous $f$ (through its labels) and attains the fan-in rate when its funnel
assignment is sufficient (through its activations).
\end{conjecture}

\begin{remark}[When should a forest beat a single tree?]\label{rem:crossover}
Comparing the upper bounds suggests where to look, without proving a separation.  A
noise-free target lets a single tree interpolate, with interior error of order
$N^{-2/n}$ for a $C^2$ target; each tree of a funnel forest regresses an outcome that is
noisy given its own cue (the coordinates it does not perceive act as noise, of variance
$\Var(Y\mid X_{R(\tau)})$), with error of order $N^{-2/(k+4)}$.  The bounds favour the
forest only when $n>k+4$.  For a noisy target the single tree's bound is
$N^{-2/(n+4)}$ and the comparison favours the forest as soon as $n>k$.
\end{remark}

%=====================================================================
\section{Inferring tasks}\label{sec:infer}
%=====================================================================

\Cref{cor:iff} reduces learnability to choosing a sufficient task assignment.  The
assignment is part of the forest's structure; we ask whether it can be inferred from the
stream, and by what rule.

\subsection{Tasks cannot be inferred locally}

\begin{definition}[Local task rule]\label{def:local}
A \emph{local task rule} decides something about a tree $\tau$---keep or prune it,
which task to give it, which message to pass---as a measurable function of $\tau$'s own
training stream (the instances $(C_{\tau,i},V_{\tau,i})$, $i\le N$, for a candidate
outcome $V_\tau$, possibly $\varnothing$) and of independent randomness.  The tree's own
prequential score, its category utility and its learned concepts are all functions of
this stream.
\end{definition}

\begin{theorem}[Locality obstruction]\label{thm:locality}
Let $X_1,X_2,X_3$ be independent and uniform on $[-1,1]$, observe only $(X_1,X_2,Y)$,
and consider the two worlds
\[
  P_a:\ Y=\operatorname{sign}(X_1X_2),\qquad P_b:\ Y=\operatorname{sign}(X_1X_3).
\]
(i) For $j\in\{1,2\}$, the pair $(X_j,Y)$ has the same law under $P_a$ and $P_b$: $X_j$
and $Y$ are independent and $Y$ is uniform on $\{-1,1\}$.  Hence, for every candidate
outcome, the training stream of a tree with cue $X_j$ has the same law in both worlds,
and so does the output of every local task rule.  (ii) Under $P_a$, $f(X)=Y$ and
discarding $X_2$ loses everything ($\E\Var(Y\mid X_1)=1$ while $\E\Var(Y\mid X)=0$);
under $P_b$, $X_2$ is irrelevant ($f(X)=\E[Y\mid X_1]=0$).  Consequently every local rule that decides whether to keep the tree perceiving $X_2$ errs with probability at
least $1/4$ in each of the two worlds.  (iii) Under $P_a$ the funnel activations of both
univariate trees vanish, so every funnel forest built on them has the maximal deficit
$\Delta=\Var(Y)=1$, whereas the label forest of \cref{thm:labelcons} and a single tree
over $(X_1,X_2)$ are sufficient.
\end{theorem}
\begin{proof}
(i) Under $P_a$, $Y=\operatorname{sign}(X_1)\operatorname{sign}(X_2)$; given $X_2$, the
factor $\operatorname{sign}(X_1)$ is uniform on $\{-1,1\}$ and independent of $X_2$, so
$Y$ is uniform and independent of $X_2$; by symmetry the same holds for $X_1$.  Under
$P_b$, $Y=\operatorname{sign}(X_1)\operatorname{sign}(X_3)$ is independent of $X_2$
outright and, by the same argument, of $X_1$.  The stream of a tree with cue $X_j$ is an
i.i.d.\ sample of $(X_j,V)$ with $V\in\{Y,\varnothing\}$ (or another observable whose
joint law with $X_j$ is the same in both worlds, such as $X_{3-j}$), so it has the same
law in both worlds, and so does any measurable function of it.  (ii) Direct.  If $p$ is
the probability that the rule keeps the tree, it errs with probability $1-p$ under
$P_a$ and $p$ under $P_b$.  (iii) $\E[Y\mid X_j]=0$ for $j=1,2$ by (i), so
\cref{prop:xy}'s argument applies; the label and single-tree statements follow from
\cref{thm:labelcons,thm:consistency}.
\end{proof}

\begin{corollary}[Parity requires composition]\label{cor:parity}
Let $X_1,\dots,X_n$ be independent and uniform on $[-1,1]$ and
$Y=\prod_{j\in J}\operatorname{sign}(X_j)$.  For every $S$ that does not contain $J$,
$Y$ is independent of $X_S$.  Hence no tree (and no rule using the streams of any set of
trees) whose joint cue misses a coordinate of $J$ can detect the dependence, every
funnel activation of such a tree vanishes, and a forest learns $Y$ only if some tree
composes---by labels or by a merged cue---information from all of $J$.
\end{corollary}
\begin{proof}
For $j\in J\setminus S$, $\operatorname{sign}(X_j)$ is uniform on $\{-1,1\}$ and
independent of $X_S$ and of the other factors, so $Y$ is uniform and independent of $X_S$.
\end{proof}

\Cref{thm:locality} is not a peculiarity of Cobweb: synergistic information---information
about $Y$ held by a set of variables jointly and by none of them alone
\citep{williams2010}---is invisible to any module that sees one of them.  Back-propagation
escapes it because each module receives the error of the composition.  A Cobweb tree
receives no error, so the decision has to be taken where the information is composed.

\subsection{Verification at the point of composition}

\begin{theorem}[Task selection by verification]\label{thm:selection}
Let $\cA$ be a finite set of task assignments.  For each $A\in\cA$ let $\hat F_{A,N}$ be
the output of the forest that Cobweb learns with assignment $A$ from a training stream
of length $N$, with $|\hat F_{A,N}|\le B$ and $\E\int|\hat F_{A,N}-F_A|^2d\mu\to0$
(\cref{thm:ideal,thm:labelcons,thm:consistency} give this for activation forests, label
forests and single trees).  Let $\hat A_N$ minimise the mean squared error of $Y$ over
the training stream itself (no held-out data are needed).  Then
\[
  \Pr\Bigl[\Delta(\hat A_N)=\min_{A\in\cA}\Delta(A)\Bigr]\to1.
\]
In particular, if $\cA$ contains a sufficient assignment, the selected assignment is
sufficient with probability tending to one and the selected forest is consistent.
\end{theorem}
\begin{proof}
Let $\sigma^2=\E\Var(Y\mid X)$.  Conditionally on the training stream, the expected
validation error of $A$ is
$R_N(A)=\sigma^2+\E[(f-\hat F_{A,N})^2\mid\text{train}]$, and
$\E[(f-\hat F_{A,N})^2\mid\text{train}]\to\Delta(A)$ in probability, because
$\|f-\hat F_{A,N}\|_{L_2(\mu)}-\|f-F_A\|_{L_2(\mu)}$ is bounded by
$\|\hat F_{A,N}-F_A\|_{L_2(\mu)}\to0$.  The validation average of the bounded squared
errors differs from $R_N(A)$ by $O_P(N_v^{-1/2})$ (Chebyshev).  With
$\gamma=\min\{\Delta(A)-\Delta^*:\Delta(A)>\Delta^*\}>0$, $\Delta^*=\min_A\Delta(A)$,
every assignment with $\Delta(A)>\Delta^*$ has a larger validation error than every
minimiser with probability tending to one, and $\cA$ is finite.
\end{proof}

The validation error at the output is the outcome of the composing tree.  With the log
loss in place of the squared error it estimates the conditional entropy
$H(Y\mid U_o)$, so \cref{thm:selection} selects the task assignment that maximises the
predictive information $I(Y;U_o)$---the same quantity, mutual information, that category
utility maximises when Cobweb chooses among its operators, evaluated one level up.  In
the online protocol the prequential error $\frac1N\sum_i(Y_i-\hat F_{A,i-1}(X_i))^2$
plays the role of the validation error \citep{dawid1984}; it converges to the same limit
whenever the online forest's risk tends to that of $F_A$, by a Ces\`aro argument.

\subsection{Task operators}

\begin{definition}[Task operators and synergy deficit]\label{def:operators}
Let a tree with outcome $V$ compose two subtrees $\sigma_1,\sigma_2$ with cues
$C_1,C_2$.  The \emph{task operators} at this composition are
\begin{itemize}[leftmargin=*,noitemsep]
\item \textsc{act}: $\sigma_j$ gets the outcome $V$ and passes its activation;
\item \textsc{label}: $\sigma_j$ gets no outcome and passes its label;
\item \textsc{hybrid}: $\sigma_j$ gets the outcome $V$ and passes its activation and label;
\item \textsc{cross}: $\sigma_j$ gets a coordinate of $C_{3-j}$ as outcome and passes its
  activation;
\item \textsc{merge}: $\sigma_1,\sigma_2$ are replaced by one tree with cue $C_1\cup C_2$.
\end{itemize}
The \emph{synergy deficit} of the pair is
$\Delta_{\textsc{act}}(\sigma_1,\sigma_2)=\E\Var(V\mid a_1,a_2)-\E\Var(V\mid C_1,C_2)\ge0$,
with $a_j=\E[V\mid C_j]$; it is zero exactly when \textsc{act} loses nothing at this
composition, and it is estimated by the difference of the validation errors of
\textsc{act} and \textsc{merge}.
\end{definition}

\textsc{merge} and its inverse mirror Cobweb's own merge and split one level up: they
decide whether two cues are perceived by one concept hierarchy or by two.

\begin{procedure}[Task induction by verification]\label{proc:induction}
\leavevmode
\begin{enumerate}[leftmargin=*,noitemsep]
\item Start with one tree per raw coordinate (or per given block).
\item Estimate the synergy deficit of every pair; merge the pairs whose deficit is
  significantly positive (union--find over these pairs).
\item Compose the resulting groups bottom-up with two-input trees; at each composition
  choose the operator in $\{\textsc{act},\textsc{label},\textsc{hybrid},\textsc{cross}\}$
  by the validation error of the composing tree's outcome.
\end{enumerate}
\end{procedure}

By \cref{thm:selection} each choice in step~3 is consistent given the subtrees below it;
the procedure as a whole is greedy and carries no global guarantee.  By
\cref{thm:locality} no step can be replaced by a rule that looks at a subtree alone.
Every tree still learns only with \texttt{ifit}: the procedure chooses \emph{what} a tree
predicts---which observable it will see---never how it learns.

%=====================================================================
\section{Experiments}\label{sec:experiments}
%=====================================================================

\emph{(Being filled in from \texttt{results/}.)}

\begin{center}
\small
\begin{tabular}{@{}lll@{}}\toprule
Claim & Test & Prediction \\\midrule
\Cref{thm:single}, \eqref{eq:singlerate} & $\sin(\pi x)$, $J_0(20x)$, one tree & error $\propto N^{-1/d}$ (up to logs) \\
\Cref{prop:prior} & same, four weightings & \textsc{prior}, \textsc{posterior} (all concepts) plateau \\
\Cref{rem:rates,rem:crossover} & 4-D and 100-D KAN functions & forest vs.\ single tree by $n$, $k$, noise \\
\Cref{prop:suffclass}, \cref{cor:iff} & $\exp(\sin\pi x+y^2)$, 4-D, 100-D & funnel forest consistent \\
\Cref{prop:xy}, \cref{thm:locality} & $xy$, $\operatorname{sign}(x_1x_2)$ & funnel error $=\mathrm{sd}(f)$; labels repair \\
\Cref{thm:locality} & sign worlds $P_a$, $P_b$ & local scores identical; composing score separates \\
\Cref{thm:selection}, \cref{proc:induction} & selection frequencies; synergy matrices & sufficient task chosen; groups recovered \\
\Cref{prop:multiview} & Gaussian two-view model & cross-view forest consistent without seeing $Y$ \\
\Cref{thm:labelcons} & label cells & cells shrink; error $\downarrow$, slowly \\
\Cref{lem:leaves} & continual 5-peak task & no forgetting \\
--- & Feynman, housing, MNIST, text & forest vs.\ single tree vs.\ MLP/KAN \\\bottomrule
\end{tabular}
\end{center}

\appendix
%=====================================================================
\section{Verification}\label{app:verify}
%=====================================================================

\paragraph{Executable verifier.}
\texttt{verification/checks.py} instantiates the objects each statement talks about with
the real C++ tree---learned by \texttt{ifit}, or built exactly for the constructions in
the proofs (depth-one grid trees, the label-forest construction, interwoven trees)---and
tests the statement's conclusion; \texttt{verification/run.py} writes the table below.
The first check compares \cref{def:predict} and \cref{def:concept}, implemented from the
definitions, with the code's scores and mixtures.

\IfFileExists{tables/verification.tex}{\begin{center}\small\input{tables/verification}\end{center}}{}

\paragraph{Adversarial review.}
Each statement is also checked by an independent reviewer instructed to break it; the
reviewer is calibrated beforehand on a copy of this document with planted errors.  The
protocol and its results are in \texttt{verification/REVIEW.md}.

\bibliographystyle{plainnat}
\bibliography{refs}
\end{document}
